\documentclass[11pt,reqno]{amsart}

%==============================================================
% Curvature on the Hurst line:
% graded Bakry--Emery curvature, Wiener chaos,
% and iterated Brownian motion
%
% Sections 2-6 are proved or numerically certified; Section 7 is a
% deliberately speculative ladder, written to be peeled back; Section 8
% is the ledger separating the two.
%==============================================================

\usepackage[utf8]{inputenc}
\usepackage[T1]{fontenc}
\usepackage{lmodern}
\usepackage[margin=1.15in]{geometry}
\usepackage{amsmath,amssymb,amsthm,mathtools}
\usepackage{microtype}
\usepackage{enumitem}
\usepackage{hyperref}
\emergencystretch=3em

\theoremstyle{plain}
\newtheorem{theorem}{Theorem}[section]
\newtheorem{proposition}[theorem]{Proposition}
\newtheorem{lemma}[theorem]{Lemma}
\newtheorem{corollary}[theorem]{Corollary}
\newtheorem{conjecture}[theorem]{Conjecture}
\theoremstyle{definition}
\newtheorem{definition}[theorem]{Definition}
\newtheorem{speculation}{Speculation}
\theoremstyle{remark}
\newtheorem{remark}[theorem]{Remark}

\newcommand{\R}{\mathbb R}
\newcommand{\C}{\mathbb C}
\newcommand{\N}{\mathbb N}
\newcommand{\Z}{\mathbb Z}
\newcommand{\E}{\mathbb E}
\newcommand{\T}{\mathbb T}
\newcommand{\LP}{\mathcal{LP}}
\newcommand{\Had}{\circ}
\DeclareMathOperator{\tr}{tr}
\DeclareMathOperator{\sgn}{sgn}
\DeclareMathOperator{\rank}{rank}

\title[Curvature on the Hurst line]{Curvature on the Hurst Line:\\
Graded Bakry--\'Emery Curvature, Wiener Chaos,\\
and Iterated Brownian Motion}
\author{Ramiro Fontes}
\address{Quijotic Research}
\email{ramirofontes@gmail.com}
\urladdr{www.quijotic.com}
\date{\today}

\begin{document}

\begin{abstract}
We propose that three apparently separate evolutions---the
fractional-Laplacian (Cauchy) Fokker--Planck semigroup, the de
Bruijn--Newman heat flow on the Riemann $\Xi$-function, and the
iterated-Brownian / biharmonic pseudo-process $e^{t\partial_x^4}$---are
sections of one object: the analytic continuation in the Hurst
parameter $H=\gamma/2$ of a \emph{graded} Bakry--\'Emery curvature
hierarchy. The continuation runs along a single line $H\in(0,\infty)$:
genuine jump processes and fractional Brownian motion for $H\in(0,1)$;
the self-dual heat point at $H=1$; and signed pseudo-processes,
iterated Brownian motion among them, for $H>1$.

The hard content is a closed form for the whole hierarchy at the
Cauchy point. Writing $R=R_{1/2}=[\min(n,m)]$ for the standard
Brownian covariance and $R^{\Had j}$ for its $j$-th Hadamard power
(the kernel of the $j$-th iterated carr\'e du champ $\Gamma_j$), we
prove that $R^{-1}R^{\Had j}$ is upper triangular with spectrum
$\{\,k^{j}-(k-1)^{j}\,\}_{k=1}^N$: the discrete derivative of the
power $k\mapsto k^j$. This recovers the odd integers $\{1,3,\dots,2N-1\}$
at $j=2$, gives $\tr=N^j$, $\det=\prod_k(k^j-(k-1)^j)$, and a uniform
\emph{higher} curvature bound $\Gamma_j\ge\Gamma$ at every chaos level
$j$, since the minimum eigenvalue is $1$ for all $j$. The clean engine
is a Cholesky factorization of the chaos covariance,
$R^{\Had j}=LD_jL^{\top}$ with $L$ the summation operator and
$D_j=\mathrm{diag}(k^j-(k-1)^j)$ its innovation variances; the whole
tower commutes, is generated by a single number operator
$\mathcal N$ via $M^{(j)}=(\mathcal N+1)^j-\mathcal N^j$, and the
universal bound $\kappa_j=1$ is the curvature of the Fock vacuum.
Its secular polynomial $z^j-(z-1)^j$ has all zeros on the critical line
$\Re=\tfrac12$, so the whole tower is totally nonnegative and its
positivity is a provable ``critical line $\Rightarrow$ positivity''---a
solvable model of the mechanism the Riemann hypothesis would need.
The grading is
the Wiener chaos of the Brownian dual, $R^{\Had j}=\tfrac1{j!}\,
\mathrm{Cov}(j\text{-th chaos})$. Numerically, $\gamma=1$ maximizes the
$j$-th curvature for every $j$, extending the known $j=2$ phase
transition to the whole tower.

At the iterated-Brownian endpoint $H=2$ ($\gamma=4$) the base kernel
$\Psi_4$ is a rank-three form of signature $(1,2)$, the curvature shadow
of the Laguerre--P\'olya non-preservation of $e^{t\partial_x^4}$. Where
curvature meets real-rootedness we make precise through a
total-positivity dictionary: among even orders, the kernel is positive
$\Leftrightarrow$ the flow preserves $\LP$ $\Leftrightarrow$ heat, and the
curvature spectral zeta at $j=2$ is the Riemann zeta itself,
$\sum_k(2k-1)^{-s}=(1-2^{-s})\zeta(s)$. The signature also drives the quartic
de Bruijn--Newman constant $\Lambda(4)$, which we resolve by GPU
computation: the re-coherence time of a drowned pair is independent of both
gap and height (measured exponent $\beta\approx0$ over a decade of gap and
$x$ to $6000$), because it is set by the detachment of the saddle of
$g(u)=tu^4+\log\Phi(u)$ at $u_*=\tfrac34$, $t_*=\tfrac{64}{27}\pi e^3\approx
150$---a phase with no $x$ in it, so height-blind. Every drowned pair, to
height $x\approx10^4$ far past the window, is extinguished at this scale,
not migrated outward; a direct scan confirms the bad set empty past
$t\approx160$. Thus $\Lambda(4)\approx155$ is \emph{finite} for the
operative mechanism both source papers debate; even the deep far field at
$\Phi$'s strip edge---the IBM open estimate, below float64---is reached in
$460$-digit arithmetic and found real-rooted at $t=300$, leaving only
uniformity in $t$ unproved. The same saddle gives $\Lambda(2m)<\infty$ at
every even order.
\end{abstract}

\maketitle

%==============================================================
\section{Introduction}\label{sec:intro}
%==============================================================

\subsection{One line, three flows}
Fix the symbol family $|\xi|^\gamma$ and read it twice. Descending,
$e^{-t|\xi|^\gamma}$ is the $\gamma$-stable Markov semigroup; on the
torus $\T=\R/2\pi\Z$ its Bakry--\'Emery curvature is controlled, by the
stable--fBM correspondence~\cite{FontesBE}, by the covariance of
fractional Brownian motion with Hurst parameter $H=\gamma/2$.
Ascending, $e^{+t\,u^\gamma}$ is the de Bruijn--Newman multiplier; at
$\gamma=2$ it is P\'olya's universal factor that drives the zeros of
$\Xi$ to the line, and at $\gamma=4$ it is the iterated-Brownian /
biharmonic pseudo-process $e^{t\partial_x^4}$ of Funaki and
Allouba--Zheng, which does \emph{not} preserve the Laguerre--P\'olya
class~\cite{FontesIBM}.

These are usually kept in separate worlds: a jump diffusion, an
arithmetic flow, a signed pseudo-process. We organize them on a single
axis, the \emph{Hurst line} $H=\gamma/2\in(0,\infty)$:
\[
\underbrace{H\in(0,1)}_{\text{jump / fBM / curvature}}
\ \Big|\
\underbrace{H=1}_{\text{heat: self-dual}}
\ \Big|\
\underbrace{H>1}_{\text{signed pseudo-process}} .
\]
The wall at $H=1$ ($\gamma=2$) is the Schoenberg threshold: $|\xi|^\gamma$
is conditionally negative definite, equivalently fBM with Hurst $H$
exists, equivalently the carr\'e du champ kernel is positive
semidefinite, \emph{iff} $H\le 1$~\cite{Schoenberg}. To the left lives
the curvature theory of~\cite{FontesBE}, with its champion the Cauchy
point $H=\tfrac12$. At the wall lives the heat semigroup, the one point
that is simultaneously a Markov diffusion (classical Bakry--\'Emery)
and a Laguerre--P\'olya preserver (de Bruijn--Newman): the self-dual
point. To the right, past the wall, the covariance loses positivity,
the process becomes a signed pseudo-process, and the first integer
station is $H=2$: iterated Brownian motion.

\subsection{What is actually proved}
The spine of the paper is one computation. The Bakry--\'Emery hierarchy
of~\cite{FontesBE} attaches to the $j$-th iterated carr\'e du champ
$\Gamma_j$ the kernel $\Psi_\gamma^{\Had j}$ (the $j$-th Hadamard
power), which on same-sign frequencies is $R_H^{\Had j}$. At the
Cauchy point $H=\tfrac12$, where $R_{1/2}(n,m)=\min(n,m)$ is the
covariance of standard Brownian motion, we evaluate the whole tower in
closed form.

\begin{theorem}[Graded curvature spectrum; \S\ref{sec:graded}]
\label{thm:intro-graded}
Let $R=[\min(n,m)]_{n,m=1}^N$. For every integer $j\ge1$ the matrix
$M^{(j)}:=R^{-1}R^{\Had j}$ is upper triangular, with spectrum
\[
   \mathrm{spec}\,M^{(j)}=\bigl\{\,k^{j}-(k-1)^{j}\ :\ k=1,\dots,N\,\bigr\}.
\]
Consequently $\min\mathrm{spec}\,M^{(j)}=1$ for all $j$ (the higher
curvature bound $\Gamma_j\ge\Gamma$ holds at every chaos level),
$\tr M^{(j)}=N^{j}$, and $\det M^{(j)}=\prod_{k=1}^N\bigl(k^j-(k-1)^j\bigr)$,
which at $j=2$ is $(2N-1)!!=\E[Z^{2N}]$, $Z\sim\mathcal N(0,1)$.
\end{theorem}

The numbers $k^j-(k-1)^j$ are the discrete derivative of $k\mapsto k^j$;
the theorem says the chaos-graded curvature spectrum is precisely the
discretization of $\tfrac{d}{dk}k^{j}=jk^{j-1}$. The case $j=2$ is the
odd-integer theorem of~\cite{FontesBE}; the present statement promotes
it from a single curiosity to a graded law, and identifies the grading:
$R^{\Had j}$ is, up to $j!$, the covariance of the $j$-th Wiener chaos
of the Brownian dual (\S\ref{sec:graded}). The Bakry--\'Emery tower is
the Wiener chaos decomposition seen through the carr\'e du champ.

Two further facts pin the endpoints of the Hurst line.

\begin{proposition}[All-orders phase transition; \S\ref{sec:phase}]
\label{prop:intro-phase}
For $j=2,3,4$ and $N$ up to $40$, the curvature
$\kappa_j(\gamma)=\lambda_{\min}(R_{\gamma/2}^{\Had j},R_{\gamma/2})$ is
maximized over $\gamma\in(0,2)$ uniquely at $\gamma=1$, with
$\kappa_j(1)=1$. The Cauchy point is the curvature champion at every
chaos level, not only at $j=2$.
\end{proposition}

\begin{proposition}[Collapse at the iterated-Brownian endpoint;
\S\ref{sec:ibm}]\label{prop:intro-ibm}
At $H=2$ ($\gamma=4$) the carr\'e du champ kernel
$\Psi_4(\xi,\eta)=\tfrac12(\xi^4+\eta^4-(\xi-\eta)^4)$ of the biharmonic
generator $\partial_x^4$ has rank $3$ and signature $(1,2)$: a single
positive direction against two negative ones. The diffusion's positive
cone has collapsed to a line.
\end{proposition}

The factorization behind the spectrum keeps giving. The chaos covariance
Cholesky-factors as $R^{\Had j}=LD_jL^{\top}$ with $L$ the summation
operator and $D_j=\mathrm{diag}(k^j-(k-1)^j)$ (\S\ref{sec:innov}); from
this the whole tower commutes and is the functional calculus of a single
number operator, $M^{(j)}=(\mathcal N+1)^j-\mathcal N^j$
(Prop.~\ref{prop:numberop}); is totally nonnegative, all minors given by
Cauchy--Binet (Prop.~\ref{prop:TN}); and has secular polynomial
$z^j-(z-1)^j$ with all zeros on the critical line $\Re=\tfrac12$
(Prop.~\ref{prop:critline}), so the tower's positivity is a provable
``critical line $\Rightarrow$ positivity.'' Each of these is the Cauchy
point alone---a sixth and seventh fingerprint to add to the four
of~\cite{FontesBE}. At the classical level $j=2$ the curvature spectral
zeta is the Riemann zeta itself, $\sum_k(2k-1)^{-s}=(1-2^{-s})\zeta(s)$
(Prop.~\ref{prop:czeta}). Finally, where the curvature meets
real-rootedness is made precise by the total-positivity dictionary of
\S\ref{sec:dict}, whose even-order core---kernel positive
$\Leftrightarrow$ flow preserves $\LP$ $\Leftrightarrow$ heat---is
Theorem~\ref{thm:dict}.

\subsection{What is speculation}
The proved spine above grew out of a speculation that mostly did not
survive contact with computation, and \S\ref{sec:spec} keeps the ladder
of statements $S1,\dots,S5$ precisely so the reader can see where it gave
way. Two rungs fell and left theorems behind: the ``two Schur theorems''
of $S1$ became the dictionary (\S\ref{sec:dict}), and the load-bearing
$S2$---``tower positivity \emph{is} real-rootedness''---turned out
\emph{false} (tower positivity is test-function-independent) and was
replaced by total positivity (\S\ref{sec:dict}) and the critical-line
reading (\S\ref{sec:transport}). One numerical prediction (the $\Lambda(4)$
re-coherence law of \S\ref{sec:lambda4}) was tested and refuted in its
first form. What remains genuinely speculative is $S3$--$S5$: RH as
critical tangency, zero dynamics as a Ricci flow, and a single analytic
operator across the wall. Section~\ref{sec:ledger} is the ledger
separating proved, conjectured, and dreamt.

%==============================================================
\section{The graded curvature spectrum}\label{sec:graded}
%==============================================================

We recall the minimal apparatus of~\cite{FontesBE} and then compute.

\subsection{The Hadamard-power tower}
On $\T$ with $\mu=dx/2\pi$, the fractional Laplacian
$L_\gamma=-(-\Delta)^{\gamma/2}$ acts by $L_\gamma e^{inx}=-|n|^\gamma
e^{inx}$. Its sesquilinear carr\'e du champ has Fourier kernel
$\Psi_\gamma(\xi,\eta)=\tfrac12(|\xi|^\gamma+|\eta|^\gamma-|\xi-\eta|^\gamma)$,
and on same-sign frequencies $\Psi_\gamma(\xi,\eta)=R_{H}(|\xi|,|\eta|)$
with $H=\gamma/2$, the covariance of fBM~\cite{FontesBE,Mishura}. The
iterated carr\'e du champ $\Gamma_2$ has kernel
$\Phi_{\gamma,2}=\Psi_\gamma^{\Had2}$ (the entrywise square), and more
generally the $j$-th iterate $\Gamma_j$ has kernel the $j$-th Hadamard
power $\Psi_\gamma^{\Had j}=R_H^{\Had j}$ on same-sign
frequencies~\cite{FontesBE}. On the positive-frequency polynomials
$\mathcal T_N^+=\mathrm{span}\{e^{inx}:1\le n\le N\}$, the $j$-th
curvature is the generalized eigenvalue
\begin{equation}\label{eq:kappa-j}
   \kappa_j(\gamma,N)=\lambda_{\min}\!\bigl(R_H^{\Had j},\,R_H\bigr),
   \qquad H=\gamma/2 .
\end{equation}
The classical Bakry--\'Emery curvature is $j=2$. We now solve the tower
at $H=\tfrac12$.

\subsection{The Cauchy point}
At $H=\tfrac12$, $R(n,m)=\min(n,m)$ is the Green's function of the
discrete Dirichlet Laplacian on $\{1,\dots,N\}$: its inverse
$R^{-1}$ is tridiagonal, acting on a column $v$ by
\begin{equation}\label{eq:Rinv}
   (R^{-1}v)_k=-v_{k-1}+2v_k-v_{k+1}\quad(1\le k<N,\ v_0:=0),
   \qquad (R^{-1}v)_N=-v_{N-1}+v_N .
\end{equation}
Thus $R^{-1}$ is discrete differentiation against the integration
kernel $\min(n,m)$; the theorem is that applying it to the $j$-th
Hadamard power differentiates the power profile $k\mapsto k^j$.

\begin{theorem}[Graded spectrum at $H=\tfrac12$]\label{thm:graded}
For each integer $j\ge1$ let $M^{(j)}=R^{-1}R^{\Had j}$ with
$R^{\Had j}_{nm}=\min(n,m)^j$. Then $M^{(j)}$ is upper triangular with
diagonal entries
\[
   \lambda^{(j)}_k=k^{j}-(k-1)^{j}\qquad(k=1,\dots,N),
\]
so $\mathrm{spec}\,M^{(j)}=\{k^j-(k-1)^j\}_{k=1}^N$. Moreover
\[
   \min_k\lambda^{(j)}_k=\lambda^{(j)}_1=1,\qquad
   \tr M^{(j)}=\sum_{k=1}^N\bigl(k^j-(k-1)^j\bigr)=N^{j},\qquad
   \det M^{(j)}=\prod_{k=1}^N\bigl(k^j-(k-1)^j\bigr).
\]
\end{theorem}

\begin{proof}
Fix a column index $m$ and look at row $k$ of $M^{(j)}$, i.e.\ apply
\eqref{eq:Rinv} to the column $(\min(k,m)^j)_k$.

\emph{Below the diagonal, $k>m$.} Then $\min(k-1,m)=\min(k,m)=\min(k+1,m)=m$,
so for $1\le k<N$ the entry is $-m^j+2m^j-m^j=0$, and at $k=N$ it is
$-m^j+m^j=0$. Hence $M^{(j)}_{km}=0$ for $k>m$: the matrix is upper
triangular.

\emph{On the diagonal, $k=m$.} For $k<N$,
$M^{(j)}_{kk}=-\min(k-1,k)^j+2\min(k,k)^j-\min(k+1,k)^j
=-(k-1)^j+2k^j-k^j=k^j-(k-1)^j$. For $k=N$,
$M^{(j)}_{NN}=-(N-1)^j+N^j$, the same formula. Since $M^{(j)}$ is upper
triangular, its spectrum is the diagonal.

The minimum: $k\mapsto k^j$ is convex for $j\ge1$, so its forward
differences $k^j-(k-1)^j$ increase in $k$; the least is at $k=1$, value
$1^j-0^j=1$. The trace telescopes, $\sum_{k=1}^N(k^j-(k-1)^j)=N^j-0=N^j$.
The determinant is the product of the diagonal.
\end{proof}

\begin{remark}[The odd integers and Gaussian moments]
At $j=2$, $\lambda^{(2)}_k=2k-1$ recovers the odd-integer spectrum and
$\det M^{(2)}=\prod(2k-1)=(2N-1)!!=\E[Z^{2N}]$ of~\cite{FontesBE}. The
general determinant $\prod_k(k^j-(k-1)^j)$ is the chaos-level-$j$
analogue of the Gaussian moment; the trace $N^j$ is its clean
counterpart. The minimum eigenvalue $1$ \emph{at every level} is the
statement that the Cauchy point carries a uniform higher curvature
bound $\Gamma_j\ge\Gamma$, a permanence absent for any other $H$.
\end{remark}

\subsection{The grading is Wiener chaos}
The Hadamard power is not a formal device. For standard Brownian motion
$B$, $\E[B_sB_t]=\min(s,t)$, and the Wick (Hermite) powers
$:\!B_s^j\!:$ satisfy
\begin{equation}\label{eq:chaos}
   \E\bigl[:\!B_s^j\!:\,:\!B_t^j\!:\bigr]=j!\,\bigl(\E[B_sB_t]\bigr)^{j}
   =j!\,\min(s,t)^{j},
\end{equation}
so $R^{\Had j}=\tfrac1{j!}\,\mathrm{Cov}\bigl(:\!B^j\!:\bigr)$ is, up to
$j!$, the Gram matrix of the $j$-th Wiener chaos of the Brownian dual
restricted to integer times. The Bakry--\'Emery tower
$\{\Gamma_j\}_{j\ge1}$ \emph{is} the Wiener chaos decomposition of the
fBM dual seen through the carr\'e du champ, and
Theorem~\ref{thm:graded} reads its curvature off level by level. The
eigenvector at $k=1$ (curvature exactly $1$) is the chaos ground mode;
the $k$-th carries curvature $k^j-(k-1)^j\sim jk^{j-1}$, so curvature
grows polynomially up the spectrum, with rate $j$ steepening in the
chaos level---a quantitative hypercontractivity profile.

\begin{remark}[Continuum profile]\label{rem:continuum}
With $k=xN$, $\lambda^{(j)}_{xN}\sim j(xN)^{j-1}$, so the curvature
spectrum fills $[1,\sim jN^{j-1}]$ and the geometric-mean contraction
rate is $\bigl(\prod_k\lambda^{(j)}_k\bigr)^{1/N}\sim j\,(N/e)^{\,j-1}$,
generalizing the $2N/e$ of the $j=2$ volume comparison
in~\cite{FontesBE}.
\end{remark}

\subsection{An innovation factorization}\label{sec:innov}
The telescoping behind Theorem~\ref{thm:graded} is itself a matrix
factorization, and the cleanest proof comes from naming it. Let $L$ be
the $N\times N$ unit lower-triangular all-ones matrix, $L_{nk}=\mathbf
1_{k\le n}$ (discrete summation), so that $L^{-1}$ is the first-difference
$(L^{-1}v)_n=v_n-v_{n-1}$.

\begin{lemma}[Cholesky of the chaos covariance]\label{lem:chol}
For every $j\ge1$, with $D_j=\mathrm{diag}\bigl(k^j-(k-1)^j\bigr)_{k=1}^N$,
\[
   R^{\Had j}=L\,D_j\,L^{\top}.
\]
In particular $R=LL^{\top}$ (the case $j=1$, $D_1=I$).
\end{lemma}

\begin{proof}
Telescoping the power,
$\min(n,m)^j=\sum_{k=1}^{\min(n,m)}\bigl(k^j-(k-1)^j\bigr)
=\sum_{k=1}^N\mathbf 1_{k\le n}\,\mathbf 1_{k\le m}\,(k^j-(k-1)^j)
=(L\,D_j\,L^{\top})_{nm}$.
\end{proof}

\begin{corollary}[Spectrum, eigenvectors, and a commuting tower]
\label{cor:innov}
$M^{(j)}=R^{-1}R^{\Had j}=L^{-\top}D_jL^{\top}$ is similar to $D_j$;
this reproves Theorem~\ref{thm:graded}. The curvature eigenvectors are
the columns of $L^{-\top}$, the discrete differences $e_k-e_{k-1}$,
\emph{independent of the chaos level $j$}. Hence the whole tower
$\{M^{(j)}\}_{j\ge1}$ is a commuting family, simultaneously diagonalized
by the difference basis, and $\det M^{(j)}=\det D_j=\prod_k(k^j-(k-1)^j)$,
$\tr M^{(j)}=\tr D_j=N^j$ are immediate.
\end{corollary}

Because the tower commutes, it is the functional calculus of a single
operator, and that operator is recognizable.

\begin{proposition}[A number operator generates the tower]\label{prop:numberop}
Let $\mathcal N=\tfrac12\bigl(R^{-1}R^{\Had2}-I\bigr)
=L^{-\top}\,\mathrm{diag}(k-1)_{k=1}^N\,L^{\top}$, the operator with Fock
spectrum $\{0,1,\dots,N-1\}$; the Bakry--\'Emery curvature operator is
$R^{-1}R^{\Had2}=2\mathcal N+1$. Then the entire Cauchy-point tower is
\[
   M^{(j)}=(\mathcal N+1)^{j}-\mathcal N^{j}\qquad(\text{all real }j>0),
\]
the discrete derivative of $\mathcal N\mapsto\mathcal N^{j}$. The
hierarchy thus carries two gradings---frequency $k$ (the spectrum of
$\mathcal N+1$) and chaos level $j$ (the power)---and is generated as a
commutative algebra by the one operator $\mathcal N$.
\end{proposition}

\begin{proof}
In the difference basis $L^{\top}$, $\mathcal N=\mathrm{diag}(k-1)$ and
$M^{(j)}=\mathrm{diag}(k^j-(k-1)^j)$ (Corollary~\ref{cor:innov}); the
identity $k^j-(k-1)^j=(\,(k-1)+1\,)^j-(k-1)^j$ is the claim eigenvalue by
eigenvalue, hence holds as operators since both sides are functions of
$\mathcal N$.
\end{proof}

The universal bound becomes a vacuum statement: at the Fock ground state
$\mathcal N=0$, $(\mathcal N+1)^j-\mathcal N^j=1^j-0^j=1$, so
$\kappa_j=1$ for every $j\ge1$ (Theorem~\ref{thm:graded}) is exactly
\emph{the curvature of the vacuum at every chaos level}. The
factorization has, in parallel, a probabilistic reading that names the
eigenvalues.
A centered Gaussian process with covariance $R^{\Had j}$ (the Gaussian
representative of the $j$-th Wiener chaos, \eqref{eq:chaos}) is, by
Lemma~\ref{lem:chol}, the partial-sum process
\[
   X^{(j)}_n=\sum_{k\le n}\sqrt{k^j-(k-1)^j}\;\xi_k,\qquad
   \xi_k\ \text{i.i.d.}\ \mathcal N(0,1),
\]
with \emph{independent increments} of variance $k^j-(k-1)^j$. The
curvature eigenvalue at level $k$ is exactly the $k$-th innovation
(one-step prediction) variance; the chaos grading rescales the $k$-th
innovation from $1$ at $j=1$ to $k^j-(k-1)^j\sim jk^{j-1}$.

\begin{proposition}[A fifth fingerprint of the Cauchy point]
\label{prop:fifth}
Among $H\in(0,1)$, the value $H=\tfrac12$ is the unique one at which the
all-ones innovation factorization $R_H^{\Had j}=L\,D\,L^{\top}$ holds: it
already fails at $j=1$ for $H\neq\tfrac12$, since $R_H=LL^{\top}$ forces
$R_H(n,m)=\min(n,m)$, i.e.\ independent fBM increments, i.e.\ $H=\tfrac12$.
Equivalently, $L^{-1}R_H^{\Had 2}L^{-\top}$ is diagonal iff $H=\tfrac12$
(numerically: off-diagonal mass $1.5,\,0,\,10.9$ at $H=0.3,0.5,0.7$).
\end{proposition}

This adjoins to the four characterizations of the Cauchy point
in~\cite{FontesBE} (curvature maximizer; cross-sign vanishing; $H=\tfrac12$
independent increments; unit determinant) a fifth: \emph{independent
innovations survive every Hadamard power}. It is the structural reason
the spectrum is as clean as Theorem~\ref{thm:graded} says---and, by
Proposition~\ref{prop:fifth}, the reason no other $H$ shares it.

The factorization sharpens the spectral statement to an all-minors one.

\begin{proposition}[The tower is totally nonnegative]\label{prop:TN}
At the Cauchy point, for every real $j>0$ the chaos covariance
$R^{\Had j}=\min(n,m)^j$ is \emph{totally nonnegative}: every minor is
$\ge0$. Indeed $R^{\Had j}=LD_jL^{\top}$ is a product of totally
nonnegative matrices---$L$ and $L^{\top}$ have all minors in $\{0,1\}$,
$D_j$ is a positive diagonal---so total nonnegativity follows from the
Cauchy--Binet formula, which moreover evaluates each minor explicitly:
\begin{equation}\label{eq:CB}
   \det R^{\Had j}[I\,|\,J]
   =\sum_{K}\det L[I|K]\Bigl(\textstyle\prod_{k\in K}(k^j-(k-1)^j)\Bigr)
     \det L[J|K],
\end{equation}
a nonnegative sum of products of curvature eigenvalues, with leading
principal minors $\det R^{\Had j}[\{1,\dots,p\}]=\prod_{k\le p}(k^j-(k-1)^j)$.
This is once more the Cauchy point alone: for $H\neq\tfrac12$ the matrix
$R_H^{\Had2}$ already carries strictly negative minors (numerically, of
order $10$ to $10^2$ at $H=0.3,\dots,0.7$).
\end{proposition}

Total nonnegativity, not mere positive semidefiniteness, is the
positivity class of the real-rootedness dictionary
(\S\ref{sec:dict}): the Aissen--Schoenberg--Whitney--Edrei theorem reads
real zeros off the total positivity of a Toeplitz matrix, and TN
operators are exactly the variation-diminishing ones (Schoenberg). So the
Cauchy tower does not merely sit in the positive cone---it sits in the TN
cone, the same currency the dictionary trades in, sharpening both
Theorem~\ref{thm:graded} (spectrum $\to$ all minors) and the meeting point
of \S\ref{sec:dict} (PSD $\to$ TN). The one caveat: $R^{\Had j}$ is a
Green's-function kernel, not a Toeplitz matrix, so its total nonnegativity
places it in the dictionary's class without, by itself, naming a
real-rooted function.

\begin{remark}[The tower in the chaos variable, and a wall at $j=1$]
\label{rem:chaosvar}
Lemma~\ref{lem:chol} uses only telescoping, so $R^{\Had j}=LD_jL^{\top}$
holds verbatim for \emph{every real} $j>0$, and with it the trace
identity
\[
   \tr M^{(j)}=\sum_{k=1}^N\bigl(k^j-(k-1)^j\bigr)=N^{j},
\]
the discrete fundamental theorem of calculus---the case $j=2$ being
$\sum_{k\le N}(2k-1)=N^2$, ``the sum of odds is a square,'' now read as a
curvature trace. The bound $\kappa_j=\min_k(k^j-(k-1)^j)$ is attained at
$k=1$ (value $1$) for $j\ge1$, but at $k=N$ for $0<j<1$, where it is
$N^j-(N-1)^j\sim jN^{j-1}\to0$. So the higher-curvature bound
$\Gamma_j\ge\Gamma$ holds exactly for chaos levels $j\ge1$ and degrades
below: a positivity wall at $j=1$ in the chaos variable, parallel to the
Schoenberg wall at $H=1$ in the Hurst variable. The ridge
$(H,j)=(\tfrac12,[1,\infty))$, on which $\kappa\equiv1$, is an explicit
slice of the analytic operator $\mathcal K(H,j)$ dreamt in $S5$
(\S\ref{sec:spec}): at the Cauchy point the $j$-continuation is not a
dream but the closed form $k^j-(k-1)^j$.
\end{remark}

%==============================================================
\section{The phase transition at every order}\label{sec:phase}
%==============================================================

The Cauchy point is distinguished not only by its closed form but by
optimality. Reference~\cite{FontesBE} proves that
$\gamma\mapsto\kappa_2(\gamma)$ is uniquely maximized at $\gamma=1$,
with $\kappa_2(1)=1$ (the curvature phase transition). The graded tower
suggests this is one slice of an all-orders statement.

\begin{conjecture}[All-orders optimality of the Cauchy point]
\label{conj:phase}
For every integer $j\ge2$, the limit
$\kappa_j(\gamma)=\lim_{N\to\infty}\lambda_{\min}(R_{\gamma/2}^{\Had j},
R_{\gamma/2})$ exists for $\gamma\in(0,2)$ and is uniquely maximized at
$\gamma=1$, with $\kappa_j(1)=1$ and quadratic decay
$\kappa_j(1+\varepsilon)=1-c_j\varepsilon^2+O(\varepsilon^3)$.
\end{conjecture}

The evidence (computed; \texttt{numpy}/\texttt{scipy} generalized
eigenvalues, $N=40$):
\begin{center}
\begin{tabular}{rccccc}
\hline
 & $\gamma=0.5$ & $\gamma=0.8$ & $\gamma=1.0$ & $\gamma=1.2$ & $\gamma=1.5$\\
\hline
$\kappa_2(\gamma)$ & $0.9315$ & $0.9899$ & $1.0000$ & $0.9880$ & $0.8994$\\
$\kappa_3(\gamma)$ & $0.9093$ & $0.9871$ & $1.0000$ & $0.9846$ & $0.8732$\\
$\kappa_4(\gamma)$ & $0.9118$ & $0.9878$ & $1.0000$ & $0.9854$ & $0.8759$\\
\hline
\end{tabular}
\end{center}
At each level the maximum sits exactly at $\gamma=1$ with value $1$, the
profile tightening as $j$ grows. The $j=2$ proof in~\cite{FontesBE}
runs through the Z-matrix / Perron--Frobenius structure of
$R_H^{-1}R_H^{\Had2}$ for $H\le\tfrac12$; the same structure is visible
numerically for the higher powers, and a uniform proof would give
Conjecture~\ref{conj:phase} for all $j$ at once. The mechanism is the
same at every level: independence of Brownian increments at $H=\tfrac12$
makes $R_{1/2}^{\Had j}$ exactly diagonalizable against $R_{1/2}$, and
any $H\neq\tfrac12$ introduces increment correlation that distorts the
Hadamard alignment downward.

%==============================================================
\section{The iterated-Brownian endpoint: collapse to a line}
\label{sec:ibm}
%==============================================================

Cross the wall. At $H=2$ ($\gamma=4$) the stable--fBM identity of
\cite{FontesBE} continues formally,
\[
   \Psi_4(\xi,\eta)=\tfrac12\bigl(\xi^4+\eta^4-(\xi-\eta)^4\bigr)
   =R_2(|\xi|,|\eta|),
\]
the ``$H=2$ fBM covariance''---which is the carr\'e du champ of the
biharmonic / iterated-Brownian generator $\partial_x^4$ whose
multiplier $e^{tu^4}$ drives the flow of~\cite{FontesIBM}. The identity
survives; positivity does not.

\begin{proposition}[Signature of the biharmonic carr\'e du champ]
\label{prop:signature}
For $\xi,\eta>0$,
$\Psi_4(\xi,\eta)=2\xi^3\eta-3\xi^2\eta^2+2\xi\eta^3
=\xi\eta\,(2\xi^2-3\xi\eta+2\eta^2)$, a kernel of rank $3$ spanned by
$\{\xi,\xi^2,\xi^3\}$. As a quadratic form in the coordinates
$(\xi,\xi^2,\xi^3)$ its Gram matrix is
$\bigl(\begin{smallmatrix}0&0&2\\0&-3&0\\2&0&0\end{smallmatrix}\bigr)$,
with eigenvalues $\{-3,\,2,\,-2\}$: signature $(1,2)$. Hence on any set
of $N\ge2$ distinct positive frequencies the carr\'e du champ matrix
$[\Psi_4(n,m)]$ has exactly one positive eigenvalue. Already at
$N=2$ on $\{1,2\}$,
$\Psi_4=\bigl(\begin{smallmatrix}1&8\\8&16\end{smallmatrix}\bigr)$,
$\det=-48<0$.
\end{proposition}

\begin{proof}
The expansion of $\tfrac12(\xi^4+\eta^4-(\xi-\eta)^4)$ is direct. The
form is supported on the three monomials $\xi,\xi^2,\xi^3$, so its rank
is at most $3$; the displayed Gram matrix is nonsingular, so the rank
is $3$. Its characteristic polynomial factors through the antidiagonal
$2\times2$ block (eigenvalues $\pm2$) and the central entry $-3$, giving
signature $(1,2)$. A Gram matrix on $N$ frequencies is a congruence
restriction, so by interlacing it has at most one positive eigenvalue,
and at least one (the diagonal $\Psi_4(n,n)=n^4>0$). The $N=2$ minor is
$16-64=-48$.
\end{proof}

The contrast with the diffusion is total. At $H\le1$ the carr\'e du
champ is an fBM covariance: a positive cone of full dimension, on which
the whole Hadamard tower of Theorem~\ref{thm:graded} is positive. At
$H=2$ the cone has collapsed to a \emph{single line}---the genuine
energy direction $\Psi_4(\xi,\xi)=\xi^4$---against a two-dimensional
negative part. The semigroup $e^{t\partial_x^4}$ no longer contracts an
energy; it has one stabilizing mode and two destabilizing ones. This is
the spectral form of the theorem of~\cite{FontesIBM} that
$e^{t\partial_x^4}$ does not preserve $\LP$: the loss of real-rootedness
and the loss of the positive cone are the same event at $H=2$.

\begin{remark}[Even stations on the Hurst line]
For each integer $m\ge1$ ($\gamma=2m$, $H=m$), $\Psi_{2m}$ is a polynomial
kernel of rank $2m-1$, supported on the monomials $\xi,\dots,\xi^{2m-1}$
with an anti-diagonal Gram matrix (entry $\propto\xi^a\eta^{2m-a}$). Its
signature is therefore
\[
   \bigl(\,m-1+\mathbf 1_{m\ \mathrm{odd}}\,,\ \ m-1+\mathbf 1_{m\ \mathrm{even}}\,\bigr):
\]
heat $m=1$ gives $(1,0)$ (rank one, positive); past the wall the positive
and negative parts are nearly balanced---$m=3$ gives $(3,2)$, $m=4$ gives
$(3,4)$, and so on. The iterated-Brownian order $m=2$ is the exception:
signature $(1,2)$, the \emph{unique} even order beyond heat with a single
positive direction. It is the mildest crossing of the wall---one station
past, one positive direction still standing---but that mildness is special
to it, not the start of a pattern (the non-even $H>1$ give full-rank
indefinite kernels of their own).
\end{remark}

%==============================================================
\section{The $\Lambda(4)$ dichotomy, read off the signature}
\label{sec:lambda4}
%==============================================================

The quartic de Bruijn--Newman constant
$\Lambda(4)=\inf\{T:H_{t,4}\in\LP\ \forall t\ge T\}$ for
$H_{t,4}=e^{t\partial_x^4}\,\tfrac18\Xi(\cdot/2)$ is conjectured
\emph{finite} in~\cite{FontesIBM} (the flow leaves $\LP$ for a bounded
interval of times, then returns by saddle detachment) and
\emph{infinite} in~\cite{FontesRMT} (under GUE statistics, bulk
lift-off recurs at every height). The two readings are not reconciled
there. The signature of Proposition~\ref{prop:signature} suggests a
mechanism, which yields a falsifiable prediction; we follow it through a
heuristic, a CPU test that \emph{refutes} the prediction's first form, a
gap-test that revives it, and finally a GPU computation
(\S\ref{sec:gpu}) that \emph{resolves} the dichotomy: $\Lambda(4)$ is
\emph{finite}, $\approx155$, the saddle-detachment scale
$\tfrac{64}{27}\pi e^3$, with every drowned pair extinguished at that
height-independent scale (the deep strip-edge far field staying the one
open estimate). The arc is a worked example
of the ladder of \S\ref{sec:spec} climbed, peeled back, and then settled by
direct computation rather than argument. Readers wanting the answer may
skip to \S\ref{sec:gpu}; the intervening subsections record how the wrong
variable (height) masked the right one and how the saddle finally decided
it.

\subsection{The argument}
Four steps, the first proved, the rest a heuristic chain built on it.

\emph{(1) The obstruction is scale-invariant.} $\Psi_4$ is homogeneous
of degree $4$, so its signature $(1,2)$ is identical on every frequency
band: there is no high-frequency regime in which order four recovers a
positive cone. Order two is the opposite, $\Psi_2(\xi,\eta)=\xi\eta$,
rank one and positive at all scales. \emph{Order two is safe by
structure; order four can be safe only dynamically.}

\emph{(2) The signature is the random-matrix dichotomy.} The single
positive direction of $\Psi_4$ is the edge mode that stays real
(the edge-stability $e_1^2e^{2s}>0$ of~\cite{FontesRMT}); the two
negative directions are the bulk, where lift-off occurs. Signature
$(1,2)$ \emph{is} ``edge-stable, bulk-unstable.''

\emph{(3) Instability is an off-diagonal effect, and the saddle only
damps the diagonal band.} Since $\Psi_4(n,n)=n^4>0$, a single frequency
is stable; the negative cone requires at least two interacting
frequencies. The ``return'' of~\cite{FontesIBM} concentrates the flow on
one band $u_*\sim\tfrac14\log t$, killing multi-frequency interaction
and so restoring the cone---but only within the band, i.e.\ for
$|z|\le X_t\sim t^{2/3}$ (the window of~\cite{FontesIBM}). Outside the
band the kernel is still $(1,2)$-unstable, unprotected.

\emph{(4) The competition.} A close pair of $\Xi$ at height $x$ with gap
$g$ drowns under the order-four floor at
$t_{\mathrm{drown}}(x)\sim g^2/(\log x)^2\to0$: close pairs go bad
almost at once. The window reaches height $x$ only at
$t_{\mathrm{win}}(x)\sim x^{3/2}$. Having no restoring direction
(the two negative eigendirections), a drowned bulk pair stays off the
axis from $t\approx0$ until saddle arrival $\sim x^{3/2}$. Hence at every
$t$, all heights $x>t^{2/3}$ carry drowned, unprotected pairs, and
$H_{t,4}\notin\LP$ for all $t$.

\begin{conjecture}[Signature reading of $\Lambda(4)$, as a falsifiable
prediction]\label{conj:lambda4}
If the heuristic holds, the order-four flow leaves $\LP$ at a height that
marches outward like $t^{2/3}$, so the bad set is nonempty for all $t>0$
and $\Lambda(4)=\infty$. The sharp, testable form: the re-coherence time
$t^\ast(x)$ of a bad pair at height $x$ scales as $x^{3/2}$ (saddle
arrival), not as a constant.
\end{conjecture}

\subsection{The test, run: the $x^{3/2}$ law fails}\label{sec:test}
We tested Conjecture~\ref{conj:lambda4} directly. With
$H_{t,4}(x)=\int_0^\infty e^{tu^4}\Phi(u)\cos(xu)\,du$ evaluated in
high-precision interval-free quadrature (\texttt{mpmath}, $30$ digits;
the routine reproduces the Arb-certified value $L_1[H_{8,4}](63.2083)=
-2.3966\times10^{-20}$ of~\cite{FontesIBM} to all printed digits), we
located the drowning pairs near the low critical-line zeros $2\gamma_n$
and, for each, bisected the largest time $t^\ast$ at which the Laguerre
certificate $L_1[H_{t,4}]<0$ persists at the moving critical point. The
re-coherence times:
\begin{center}
\begin{tabular}{r|ccccccccc}
\hline
$x_\ast$ & $63$ & $98$ & $116$ & $120$ & $126$ & $153$ & $176$ & $191$ & $223$\\
$t^\ast$ & $37$ & $17$ & $121$ & $120$ & $118$ & $110$ & $38$ & $97$ & $88$\\
$0.09\,x^{3/2}$ & $45$ & $87$ & $112$ & $119$ & $127$ & $170$ & $211$ & $237$ & $300$\\
\hline
\end{tabular}
\end{center}
The prediction fails. The log--log slope of $t^\ast$ against $x$ is
$0.74$ (the conjecture requires $1.5$), the correlation is $0.40$, and
decisively $t^\ast$ does \emph{not} grow: its maximum $121$ is reached at
$x\approx116$ and is not exceeded out to $x\approx223$, where a pair
re-coheres at $88$ against a predicted $\sim300$. The apparent climb of
the first five pairs ($37\to120$) was a coincidence of three adjacent
pairs near $x\approx120$ drowning as one episode, not a law. Over the
whole range the bad set is empty past $t\approx121$, exactly the
``nothing past $t\approx150$'' census of~\cite{FontesIBM}.

\begin{remark}[What the failed test does and does not settle]
\label{rem:test}
At accessible heights ($x\le223$, ordinary gaps $\sim2$--$3$) the data
are consistent with the \emph{finite} reading of~\cite{FontesIBM}: the
re-coherence time is bounded, and the simple saddle-arrival law
$t^\ast\asymp x^{3/2}$ of the signature argument is refuted. The leg that
breaks is (c) of Remark~\ref{rem:loadbearing}: a drowned bulk pair at
moderate height \emph{does} re-cohere before any saddle reaches it, so
the off-diagonal instability of $\Psi_4$ is not activated as durably as
the dimension count suggested. This does not prove $\Lambda(4)<\infty$:
the infinite scenario of~\cite{FontesRMT}, if real, is carried by
anomalously close Lehmer pairs (gap $\sim10^{-2}$ near height $\sim10^4$),
which lie far outside this range and which the present scan cannot reach.
The honest verdict is that the signature \emph{mechanism} for
$\Lambda(4)=\infty$ in its first form---scale-invariant instability
outrunning a $t^{2/3}$ window---is not supported by the height scan. But
\S\ref{sec:gaptest} shows the height was the wrong variable: at fixed
small gap the re-coherence time is governed by the gap, $t^\ast\asymp
g^{-\alpha}$, which extrapolates the other way. The question stands open;
the variable, not the verdict, was the lesson.
\end{remark}

\begin{remark}[What is load-bearing]\label{rem:loadbearing}
The argument rests on three legs of unequal strength: (a) the
scale-invariant signature $(1,2)$, \emph{proved}
(Proposition~\ref{prop:signature}); (b) a supply of close bulk pairs at
every height, which is GUE statistics for $\Xi$ (Montgomery--Odlyzko),
\emph{conjectural}; (c) that a drowned bulk pair does not self-re-cohere
before saddle arrival, \emph{strongly suggested} by the absence of a
restoring eigendirection but not proved. Leg (c) is the one the test of
\S\ref{sec:test} breaks: at moderate height drowned pairs re-cohere on
their own, well before any saddle arrives, so the dimension count
$(1,2)$ overstates how durably the instability is activated. What
survives is the weaker, structural reading---signature loss at $H=2$ is
the curvature face of $\LP$-escape---not the dynamical claim that the
escape is permanent in $t$.
\end{remark}

\subsection{The gap, not the height}\label{sec:gaptest}
The refuted prediction varied the wrong variable. Section~\ref{sec:test}
fixed the gap (ordinary pairs, $2g\sim2$--$5$) and varied the height; but
\cite{FontesRMT} locates lift-off at \emph{asymmetric small-gap} pairs,
so the controlling variable should be the gap. We therefore measured
$t^\ast$ for the smallest-gap pairs among the first $120$ zeros of $\Xi$
(gaps $2g\in[1.43,1.74]$, heights $x\in[300,463]$):
\begin{center}
\begin{tabular}{r|ccccc}
\hline
$2g$ & $1.43$ & $1.45$ & $1.47$ & $1.64$ & $1.74$\\
$t^\ast$ & $133$ & $115$ & $107$ & $105$ & $83$\\
\hline
\end{tabular}
\end{center}
Now the trend is clean and monotone: $t^\ast$ \emph{rises} as the gap
shrinks, $t^\ast\asymp(2g)^{-\alpha}$ with $\alpha\approx2.4$ over this
(narrow) range, and the smallest gap gives $t^\ast=133$, already above the
height-scan ceiling of $121$. Since GUE statistics
(Montgomery--Odlyzko) furnish arbitrarily small gaps at arbitrarily large
height, a positive-power law $t^\ast\asymp g^{-\alpha}$ extrapolates to
$t^\ast\to\infty$---reviving $\Lambda(4)=\infty$, now through the correct
variable. This does not settle it: the range is narrow, height and gap are
partly confounded, and the Lehmer regime ($2g\sim10^{-2}$) remains out of
reach. But it reverses the lean of \S\ref{sec:test}: the height was a red
herring, the gap is the knob, and in the gap the data point back toward
the infinite reading of~\cite{FontesRMT}. The honest status is open, with the
mechanism relocated from ``saddle outrunning a window'' (refuted) to
``small-gap lift-off persisting as $g^{-\alpha}$'' (supported, but
unreachable in the decisive regime).

\subsection{Resolution on the GPU: $\Lambda(4)$ is finite}\label{sec:gpu}
The competition is decided by direct computation. On a GPU we evaluate
$H_{t,4}(x)=\int_0^\infty e^{tu^4}\Phi(u)\cos(xu)\,du$ and its Laguerre
certificate in double precision by Simpson quadrature on the weight's
effective support, the weight normalized by its log-maximum to stay
sign-faithful through the ten orders of oscillatory cancellation; the
routine reproduces the Arb value of~\cite{FontesIBM} at
$(t,x)=(8,63.2083)$ to all digits. Two measurements settle the question.

\emph{The re-coherence exponent is zero.} Locating the close pairs of
$\Xi$ as the small-gap sign changes of $H_0$ (the first $\sim\!10^5$ zeros,
to height $x\approx6000$), we measured the bad-interval upper edge
$t^\ast$ for pairs spanning normalized gap $s\in[0.07,0.50]$---a decade
below the CPU's reach---and height $x\in[440,6000]$:
\[
   t^\ast\in[92,162],\quad \text{mean }140,\qquad t^\ast\asymp
   s^{-0.12}\ (R^2=0.18).
\]
The re-coherence time is \emph{independent of gap and of height}:
$\beta\approx0$, nowhere near the critical $4.5$. The CPU trend
($\beta\approx2.4$, \S\ref{sec:gaptest}) was noise in a narrow band; over a
decade of gap the law is flat, and an extension to $s=0.045$ at $x\approx
10^4$---deep toward the Lehmer regime---holds the same $t^\ast\approx140$.
By the criterion $\beta<\tfrac92$, $\Lambda(4)<\infty$.

\emph{The saturation is saddle detachment, and it is height-blind.} The
flat $t^\ast$ is fixed by the phase $g(u)=tu^4+\log\Phi(u)$, with
$\log\Phi(u)\sim-\pi e^{4u}$. Its interior critical point obeys
$t=\pi e^{4u}/u^3$, minimized at $u_*=\tfrac34$, so the saddle is born at
\begin{equation}\label{eq:tstar4}
   t_*=\tfrac{64}{27}\,\pi e^{3}=149.57\qquad(u_*=\tfrac34),
\end{equation}
and overtakes the endpoint $g(0)$---where the arithmetic lives---at the
dominance crossover $t_{**}=157.0$. \emph{The phase $g$ contains no $x$},
so the detachment scale is identical at every height; past $t_{**}$ the
endpoint is suppressed by $e^{-C(t)}$, $C(t)=g(u_*)-g(0)$ climbing from $0$
to $52$ by $t=200$ and without bound after, and $H_{t,4}$ collapses to the
saddle-dominated Gaussian--cosine, real-rooted on the bulk (the window
theorem of~\cite{FontesIBM}). The endpoint suppression is uniform in $x$,
and---decisively---every drowned arithmetic pair we can reach, to height
$x\approx10^4$ far beyond the window, is extinguished at this same scale,
not migrated outward. That is the operative exit mechanism both
\cite{FontesIBM} and~\cite{FontesRMT} debate, and it is finite.

\emph{Direct confirmation.} A dense scan of $L_1[H_{t,4}]$ over
$x\in[40,6000]$, far above the window $X_t\sim t^{2/3}$, finds the
violation count full below $t_*$ and collapsing sharply at the detachment
scale---$716$ bad pairs at $t=150$, $110$ at $t=165$, $35$ at $t=200$, then
$0$ at $t=270$ and clean through $t=1500$. The collapse is centered exactly
on $t_*\approx t_{**}\approx150$--$157$.

\begin{proposition}[$\Lambda(4)$ finite]\label{prop:lam4finite}
To double-precision detectability the bad set $\{t:H_{t,4}\notin\LP\}$ is
contained in $t\lesssim t_{**}=157$, with re-coherence governed by the
$x$-independent saddle detachment of $g(u)=tu^4+\log\Phi(u)$ at
$u_*=\tfrac34$, $t_*=\tfrac{64}{27}\pi e^3$. Hence $\Lambda(4)\approx155$ is
finite. The statement is numerics-plus-saddle: double precision resolves
the arithmetic only while $C(t)\lesssim30$ ($t\lesssim185$); beyond that
the saddle dominance, not the arithmetic, carries real-rootedness.
\end{proposition}

\emph{The near far field.} The $L_1$ certificate goes blind past
$C\sim30$, so past detachment we confirm real-rootedness by a means that
does not: the \emph{positions} of the $O(1)$ real zeros. Past $t_{**}$,
$H_{t,4}(x)\approx \sqrt{2\pi}\,w\,e^{C}e^{-x^2/2|g''|}\cos(xu_*)$, whose
zeros are real and equally spaced at $\pi/u_*$. At $t=300$ and $t=600$ the
computed zeros---out to the float64 reach $x\sim\sqrt{72\,|g''|}$,
\emph{past} the window $X_t\sim|g''|^{2/3}$---are uniform to $0.4$--$0.8\%$
with no anomalous close pair (e.g.\ $x\in[326,650]$ at $t=300$,
$x\in[913,1409]$ at $t=600$): an incipient off-axis pair would betray
itself as a missing or doubled real zero, and none appears.

\emph{The strip edge, reached.} Beyond this lies the one region
\cite{FontesIBM} leaves open: the deep far field $x\gtrsim|g''|\pi/8$,
where the complex saddle $u_z$ meets the edge $\Im u=\pi/8$ of $\Phi$'s
analyticity strip---the singularity that carries the zeros of $\zeta$.
There $H\sim e^{-146}$ at $t=300$, far below float64, but \emph{not} below
high-precision arithmetic: the cancellation is $\sim200$ digits, and a
fine-grid Simpson sum at $460$ digits resolves it (spectrally accurate
because $g$ decays smoothly at the support ends; at $t=8$ it reproduces the
Arb anchor, and at the strip edge $L_1$ is stable to six digits as the grid
is refined $45\mathrm{k}\to90\mathrm{k}\to180\mathrm{k}$). At $t=300$, \emph{at}
the strip edge $x=|g''|\pi/8=2306$ and \emph{past} it, out to
$x/|g''|=0.51$ (deep in the arithmetic-dominated region, the saddle well
outside the strip), every computed point has $L_1>0$ and the real zeros
stay regular: the arithmetic does not throw a pair off the axis. The
IBM obstruction is, at this $t$, benign.

\begin{remark}[What is and isn't proved]\label{rem:gpu-proof}
Three tiers. \emph{Resolved:} the drowned-pair exit---the mechanism both
source papers debate---is extinguished at the $x$-independent saddle scale
\eqref{eq:tstar4}, measured for pairs to $x\approx10^4$, with the scale
derived in closed form. \emph{Theorem:} the bulk $|z|\lesssim X_t$ is
real-rooted (the window theorem of~\cite{FontesIBM}), and the near far
field past it is clean to double precision. \emph{Evidenced past where float64 stops:} the deep far field at the strip
edge $\Im u_z=\pi/8$---the IBM open estimate---is below float64 but
reachable in high-precision arithmetic, and is real-rooted where probed: at
$t=300$ ($460$ digits) at and past the edge, $x/|g''|$ to $0.51$,
convergence-validated; and at $t=450$ ($613$ digits) in the deep far field
$x/|g''|\approx0.38$. \emph{Still open:} \emph{uniformity}---that this holds
for all large $t$ and as $x\to\infty$, not only at the sampled $(t,x)$. So $\Lambda(4)<\infty$ is a
theorem on the bulk, established for the operative drowned-pair mechanism,
and now directly evidenced even at the strip edge; what is not proved is
the uniform far-field bound. The answer is finite, the scale is
$\tfrac{64}{27}\pi e^3$, and the residue is the same uniformity in
$t$ and height that de Bruijn--Newman theory has always ended at.
\end{remark}

The mechanism is general. For the order-$2m$ flow the phase is
$tu^{2m}-\pi e^{4u}$, saddle at $u_*=(2m-1)/4$, detachment at
$t_*(m)=2\pi e^{2m-1}/\bigl(m\,((2m-1)/4)^{2m-1}\bigr)$---$149.6$ at order
four, $101.9$ at order six, $34.3$ at order eight---finite at every even
order, deriving the permanence conjecture of~\cite{FontesIBM}. A direct
order-six scan confirms it: the bad set of $H_{t,6}$ is full at $t=100$ and
empty at $t=120$, collapsing exactly at $t_*(3)=101.9$. Among the orders
that exit ($m\ge2$) the detachment scale is \emph{largest} at order four,
where the multiplier power $2m=4$ resonates with the kernel's intrinsic
$e^{4u}$: the quartic flow carries the longest-lived exit, which is why it,
and not its neighbours, was the hard case. So $\Lambda(4)$ sides
with~\cite{FontesIBM}: finite, $\approx155$, the saddle-detachment scale.

\begin{remark}[Static obstruction, dynamic resolution]\label{rem:static-dynamic}
The signature and the saddle play complementary roles, and conflating them
was the original error. The curvature signature is \emph{static}: $\Psi_4$
indefinite means the order-four flow \emph{can} leave $\LP$---a necessary
condition, the source of the exit. Whether the exit is \emph{permanent} is
\emph{dynamic}, fixed not by the curvature but by the analytic structure of
the de Bruijn--Newman kernel $\Phi$, through the saddle of
$g(u)=tu^4+\log\Phi(u)$. The signature heuristic of \S\ref{sec:lambda4}
mispredicted $\Lambda(4)=\infty$ precisely by reading a dynamic question
(persistence) off a static datum (the indefinite cone). The resolution
required leaving the curvature framework for the kernel's saddle---which is
why the answer, $\tfrac{64}{27}\pi e^3$, carries the $e$ and $\pi$ of
$\Phi$, not the integers of the curvature tower.
\end{remark}

%==============================================================
\section{The total-positivity dictionary}\label{sec:dict}
%==============================================================

The thread that $S2$ (\S\ref{sec:spec}) will promise---that positivity
detects real-rootedness---is real, but it does \emph{not} run through the
Bakry--\'Emery tower. The tower is positive for every test function by
the Schur product theorem (the kernel is a Hadamard square,
\eqref{eq:chaos}\,ff.), so it cannot see the zeros of any particular
$f$. The genuine dictionary runs through \emph{total positivity}, where a
precise, largely classical statement holds. We record it because it is
where $\Xi$ (hence RH) and the order-four flow (hence $\Lambda(4)$)
actually sit, and because its even-order core is provable in our terms.

\subsection{Three columns}
For a real polynomial $f(x)=\sum_{k=0}^n a_k x^k$ with $a_k\ge0$, the
following are equivalent:
\begin{enumerate}[label=(\arabic*)]
\item \emph{(roots)} $f$ has only real, non-positive zeros;
\item \emph{(total positivity)} the Toeplitz matrix $(a_{i-j})_{i,j}$
  ($a_k=0$ outside $0\le k\le n$) is totally positive---a P\'olya
  frequency sequence;
\item \emph{(multiplier)} $(a_k)$ acts as a P\'olya--Schur multiplier on
  real-rooted polynomials.
\end{enumerate}
(1)$\Leftrightarrow$(2) is the Aissen--Schoenberg--Whitney--Edrei
theorem~\cite{ASW}; the bridge to (3) is P\'olya--Schur~\cite{PolyaSchur2},
$(\gamma_k)$ a multiplier sequence iff $\sum_k\gamma_k x^k/k!\in\LP$. The
engine common to all three is the Schur product theorem---Hadamard
products preserve positivity, and total positivity is the form of
positivity that records real zeros of one sign. This is the correct
content of $S1$, and the correct $S2$: real-rootedness \emph{is} a
positivity condition, on the coefficient Toeplitz/Hankel matrices, not on
the operator's curvature tower.

\subsection{The even-order coincidence (proved)}
Our family meets the dictionary cleanly at the even orders.

\begin{theorem}[Schur $=$ P\'olya--Schur on even orders]\label{thm:dict}
For integer $m\ge1$, the following are equivalent:
\emph{(i)} $m=1$;
\emph{(ii)} the carr\'e du champ kernel $\Psi_{2m}$ is positive
semidefinite;
\emph{(iii)} $e^{tu^{2m}}$ is a P\'olya--Schur multiplier, i.e.\ the flow
$e^{(-1)^m t\partial^{2m}}$ preserves $\LP$ for every $t\ge0$.
\end{theorem}
\begin{proof}
(i)$\Leftrightarrow$(ii): $\Psi_{2m}$ is positive semidefinite iff
$|\xi|^{2m}$ is conditionally negative definite iff $2m\le2$
(Schoenberg~\cite{Schoenberg}), i.e.\ $m=1$; at $m=1$,
$\Psi_2(\xi,\eta)=\xi\eta$ is rank one, while for $m\ge2$
Proposition~\ref{prop:signature} exhibits a negative eigenvalue.
(i)$\Leftrightarrow$(iii): by P\'olya's universal-factor theorem a cosine
multiplier $e^{tw^m}$ ($w=u^2$) preserves $\LP$ iff it has the admissible
form $e^{\gamma w}\prod_k(1+\delta_k w)$ with $\gamma,\delta_k\ge0$; being
entire and zero-free it must equal $e^{\gamma w}$, forcing $m=1$, and at
$m=1$ this is P\'olya's theorem~\cite{FontesIBM}.
\end{proof}

So among the even orders the Schur side (kernel positive) and the
P\'olya--Schur side (flow preserves real-rootedness) \emph{coincide},
both selecting the heat point $m=1$. This is the rigorous unification
$S1$ gestures at: not a vague analogy of two Schur theorems, but their
literal agreement on this family.

\subsection{Where RH and $\Lambda(4)$ sit}
The two hard problems occupy the two ends of the dictionary.

\emph{RH is the membership column.} $\Xi\in\LP$ iff RH, and through
(1)$\Leftrightarrow$(2) this is the total positivity of the Hankel
matrices of $\Xi$'s Taylor coefficients---equivalently the hyperbolicity
of all Jensen polynomials $J^{d,n}_\Xi$, equivalently the higher Tur\'an
inequalities. This is not a reformulation we propose but the classical
P\'olya--Csordas--Varga program~\cite{CsordasVarga2}, on which
Griffin--Ono--Rolen--Zagier proved hyperbolicity for every $d$ and
$n\gg_d 1$~\cite{GORZ}: a proved fragment of RH living in column (2).

\emph{$\Lambda(4)$ is the multiplier column.} The order-four symbol
$e^{tu^4}$ is \emph{not} a multiplier sequence (its symbol
$e^{tz^4}\notin\LP$; Borcea--Br\"and\'en~\cite{BorceaBranden2}), which is
Theorem~\ref{thm:dict}(iii) failing at $m=2$ and, on the curvature side,
the signature $(1,2)$ of Proposition~\ref{prop:signature}. Whether
$\Lambda(4)<\infty$ or $=\infty$ then asks whether \emph{iterating} a
non-multiplier returns its output to column (1)---the question
\S\ref{sec:lambda4} could not close.

\subsection{The honest meeting point with the tower}
The Bakry--\'Emery Hadamard tower does touch this dictionary, but at one
remove: through the critical-exponent theory of Hadamard
powers~\cite{BGKP}, which asks when $A^{\Had t}$ stays positive (or
totally positive). Here the answer is as strong as it can be:
Proposition~\ref{prop:TN} shows $\min(n,m)$ is a matrix whose
\emph{entire} Hadamard-power tower stays \emph{totally nonnegative}, not
merely positive semidefinite---and TN is precisely column (2)'s
currency. So the tower and the dictionary are governed by one theory of
Hadamard-power total positivity; what is false is only the literal $S2$
identification of tower positivity with the real-rootedness of a test
function. The Cauchy point is the one place the tower enters the
dictionary's cone (Proposition~\ref{prop:TN}, Cauchy-only); whether a
Green's-function TN kernel can be transported to a Toeplitz one---and so
made to name a function---is the precise form of what $S2$ wanted.

\subsection{The transport, resolved: a critical line in the tower}
\label{sec:transport}
The kernel $R^{\Had j}$ does not transport to a P\'olya frequency
(Toeplitz) matrix: its bidiagonal multipliers $\sqrt{k^j-(k-1)^j}$ are not
translation invariant, and the negative-root class of column (2) is the
wrong target. The real-rooted certificate lives elsewhere---in the
\emph{secular polynomial} of the curvature operator, and it carries a
critical line.

\begin{proposition}[Critical-line zeros of the secular polynomial]
\label{prop:critline}
Write $M^{(j)}=p_j(\mathcal N+1)$ with secular polynomial
$p_j(z)=z^j-(z-1)^j$ (Proposition~\ref{prop:numberop}). Then $p_j$ has
degree $j-1$ and \emph{all its zeros lie on the critical line}
$\Re z=\tfrac12$, namely $z_r=\tfrac12+\tfrac{i}{2}\cot\frac{\pi r}{j}$,
$r=1,\dots,j-1$. Equivalently the rotation $z=\tfrac12+iw$ makes
$q_j(w):=p_j(\tfrac12+iw)$ hyperbolic (real-rooted), $q_j\in\LP$. The
curvature eigenvalues are its values on the number spectrum,
\[
   k^j-(k-1)^j=p_j(k)=j\prod_{r=1}^{j-1}(k-z_r),
\]
positive for every integer $k\ge1$ because the spectrum of $\mathcal N+1$
lies to the right of the zero-line and the $z_r$ occur in conjugate pairs;
the vacuum value $\kappa_j=p_j(1)=1$ is the value at the integer nearest
the line.
\end{proposition}

\begin{proof}
$p_j(z)=0$ forces $z^j=(z-1)^j$, hence $|z|=|z-1|$, the perpendicular
bisector of $0$ and $1$, i.e.\ $\Re z=\tfrac12$. Solving
$z/(z-1)=e^{2\pi i r/j}$ gives $z_r=\tfrac12+\tfrac{i}{2}\cot\frac{\pi
r}{j}$. The leading term of $z^j-(z-1)^j$ is $jz^{j-1}$, giving the
product; conjugate pairing $z_{j-r}=\bar z_r$ makes
$\prod_r(k-z_r)>0$ for real $k$. Setting $z=\tfrac12+iw$ rotates the line
onto $\R$, so $q_j$ has only real zeros.
\end{proof}

\begin{remark}[Cauchy zeros for the Cauchy point]\label{rem:cauchyzeros}
As $j\to\infty$ the secular zeros fill the critical line, and their
imaginary parts $y_r=\tfrac12\cot(\pi r/j)$---with $r/j$ equidistributed in
$(0,1)$---have limiting density $\tfrac{2}{\pi(1+4y^2)}$, a Cauchy law of
scale $\tfrac12$ (the change of variables $y=\tfrac12\cot(\pi u)$ from
$u\sim\mathrm{Unif}(0,1)$). The Cauchy point's secular zeros are thus
Cauchy-distributed on the critical line---$\gamma=1$ closing on itself once
more---the empirical law matching to $2\times10^{-4}$ by $j=5000$.
\end{remark}

This relocates the bridge, and is the honest resolution of the transport
question. There are two real-rooted classes: P\'olya frequency / negative
real roots (column (2): Toeplitz total nonnegativity, ASW), and
\emph{hyperbolic / critical-line} zeros---the class the Riemann hypothesis
inhabits. The chaos tower's secular polynomials sit in the second, not the
first: $p_j$ is a stable polynomial whose zeros lie exactly on
$\Re=\tfrac12$, a \emph{provable} instance of the phenomenon RH conjectures
for $\Xi$. So the curvature tower is a \emph{solvable model of ``critical
line $\Rightarrow$ positivity''}: the positivity and total nonnegativity of
the entire Hadamard tower (Theorem~\ref{thm:graded},
Proposition~\ref{prop:TN}) \emph{follow from} the zeros of $p_j$ lying on
the line $\Re=\tfrac12$---in particular off the real number-spectrum
$\{1,2,\dots\}$, which is what forces every $p_j(k)>0$. (The implication is
one-way: mere positivity needs only no real zeros; the critical line is the
sharp, provable form it actually takes here.) Where $\Xi$'s critical-line
hypothesis is open, the
tower's is a theorem---and it is exactly what makes the curvature positive.
The transport $S2$ sought was to the negative-root class and fails; the
transport that succeeds is to the critical-line class, and there the tower
already lives.

\subsection{The curvature spectral zeta is a Riemann zeta}\label{sec:czeta}
The secular polynomial packages the eigenvalues as a product; the
\emph{spectral zeta} packages them as a Dirichlet series, and at $j=2$ it
is not an analog of a Riemann object but a Riemann object.

\begin{proposition}[$\zeta$ as the curvature spectral zeta]\label{prop:czeta}
At the Cauchy point the Bakry--\'Emery curvature operator
$M^{(2)}=2\mathcal N+1$ has spectrum the odd integers, so its spectral zeta
is
\[
   Z_2(s):=\sum_{k\ge1}(2k-1)^{-s}=(1-2^{-s})\,\zeta(s),
\]
the Dirichlet $\lambda$-function. In the critical strip its zeros coincide
with the nontrivial zeros of $\zeta$; hence the Riemann hypothesis is
exactly the statement that the critical zeros of the curvature spectral
zeta lie on $\Re s=\tfrac12$. Its zeta-regularized determinant is
$\exp(-Z_2'(0))=\sqrt2$.
\end{proposition}

The identity itself is elementary---the odd integers carry the odd part of
$\zeta$---but the reading is not: the operator whose vacuum curvature is
$1$ (Theorem~\ref{thm:graded}) and whose spectrum is the odd integers has,
for its spectral zeta, the Riemann zeta, so RH is literally a statement
about this curvature operator. Higher levels do not stay clean:
$Z_j(s)=\sum_k(k^j-(k-1)^j)^{-s}$ is a Dirichlet series with abscissa
$1/(j-1)$ and a pole there, but it is not $j^{-s}\zeta((j-1)s)$ (the
corrections are $O(1)$, not small), and its zeros are not visibly Riemann
zeros for $j\ge3$. Only the classical curvature level $j=2$ carries $\zeta$
exactly.

Thus the Cauchy-point spectrum $\{2k-1\}$ meets the critical line twice
over: through the secular polynomials $p_j$ ($j\ge3$), whose zeros lie on
$\Re=\tfrac12$ provably (Proposition~\ref{prop:critline}), and through the
spectral zeta $Z_2=(1-2^{-s})\zeta$, whose critical zeros lie on
$\Re=\tfrac12$ iff RH. One critical line is a theorem inside the operator;
the other is the Riemann hypothesis itself.

%==============================================================
\section{A speculative ladder}\label{sec:spec}
%==============================================================

We now overreach on purpose. Each rung extends the last; the intent is
to climb until the structure gives out, then record where. Throughout,
``the tower'' means the Hadamard-power curvature hierarchy
$\{R_H^{\Had j}\}_{j\ge1}$ and ``positivity'' its membership in the
positive cone.

\subsection*{$S1$. Two Schur theorems, one phenomenon}
The Hadamard powers preserve positive semidefiniteness by the Schur
product theorem; that is exactly why the tower is positive for $H\le1$.
On the ascending side, the operation that preserves real-rootedness is
the P\'olya--Schur multiplier / Schur--Szeg\H{o} composition. We
speculate these are one statement read on two sides of the
stable--fBM correspondence: \emph{the Schur product theorem for fBM
covariances and the P\'olya--Schur theorem for $\LP$ multipliers are
the descending and ascending faces of a single positivity principle,
joined at the self-dual heat point.} This rung did not stay speculative:
\S\ref{sec:dict} makes the joint precise through total positivity, and on
the even orders the two faces provably coincide (Theorem~\ref{thm:dict}).
$S1$ is kept here only to mark where the proved dictionary came from.

\subsection*{$S2$. Positivity \emph{is} real-rootedness---peeled back}
The tempting sharpening of $S1$ was the equation ``$f\in\LP\iff$ the
curvature tower of $f$ is positive at every level.'' \emph{This is false
as stated}, and instructively so: the Bakry--\'Emery tower is positive
for every $f$ (Schur product theorem), so it cannot detect the roots of
any particular $f$; the map ``$f\mapsto$ tower'' was never defined. The
correct statement does not use the operator tower at all---it uses total
positivity of the \emph{coefficient} matrices, and it is the proved
dictionary of \S\ref{sec:dict}: real-rootedness $\Leftrightarrow$ Toeplitz
total positivity $\Leftrightarrow$ P\'olya--Schur multiplier, with the
Schur product theorem as the shared engine. This is the one rung the
ladder loses outright; what replaces it (\S\ref{sec:dict}) is sturdier
than the speculation was, including the proved even-order coincidence
(Theorem~\ref{thm:dict}) and the exact slots for RH and $\Lambda(4)$.

\subsection*{$S3$. The heat point is critical, and RH lives there}
At $H=1$ the covariance $R_1(s,t)=st$ is rank one---the positive cone is
a single ray, the boundary case where the tower is positive but
degenerate. This is the unique self-dual point: Markov diffusion and
$\LP$-preserver at once. We speculate the Riemann hypothesis is the
statement that $\Xi$ sits \emph{exactly} on this critical ray:
\begin{equation}\label{eq:S3}
   \text{RH}\quad\Longleftrightarrow\quad
   \text{the curvature tower of $\Xi$ is critically positive at }H=1,
\end{equation}
and the de Bruijn--Newman constant $\Lambda$ is the curvature deficit:
the amount of heat-point ($H=1$) flow needed to push the tower into the
strict cone, with $\Lambda\le0$ iff $\Xi$ began inside it. This recasts
$\Lambda\ge0$~\cite{RodgersTao} as ``$\Xi$ is not strictly inside the
cone'' and RH as critical tangency. \emph{Peel-back note:} with $S2$
gone, the honest version of \eqref{eq:S3} routes through the dictionary
of \S\ref{sec:dict}: RH $\Leftrightarrow\Xi\in\LP\Leftrightarrow$ total
positivity of $\Xi$'s Hankel matrices (Jensen hyperbolicity), of which
only finite fragments are proved~\cite{GORZ}. The ``critical tangency at
$H=1$'' picture is then a way of \emph{reading} that membership, not a
new route to it.

\subsection*{$S4$. Zero dynamics as Ricci flow on the log-gas}
The companion random-matrix study~\cite{FontesRMT} writes the
order-$2m$ flow as zero dynamics $\dot z_j=-p^{(2m)}(z_j)/p'(z_j)$, a
three-body law at $m=2$, and finds the edge Laguerre quantity stays
positive ($\propto e^{2s}>0$) while bulk pairs lift off. We speculate
this is the BE curvature continued to $H=m$ acting as a Ricci flow on
the configuration of zeros: the surviving positive direction of
$\Psi_4$ (signature $(1,2)$, Proposition~\ref{prop:signature}) is the
edge mode that stays real, and the two negative directions are the bulk
instability that produces the first complex pair. Section~\ref{sec:lambda4}
turned this analogy into a quantitative prediction (re-coherence time
$\asymp x^{3/2}$, hence $\Lambda(4)=\infty$) and then \emph{refuted it}:
the bulk pairs re-cohere on their own at moderate height, so the
``curvature healing'' is not confined to the saddle window as the
scale-invariant signature suggested. \emph{Peel-back note:} what remains
is the static correspondence of signs---edge positive, bulk negative,
in the $(1,2)$ signature---not the dynamical claim that the bulk
instability persists in time; \S\ref{sec:test} is where that claim was
tested and failed.

\subsection*{$S5$. One operator, a bifurcation across the wall}
The boldest rung. Theorem~\ref{thm:graded} gives a spectrum
$k^j-(k-1)^j$ that is manifestly analytic in $j$ and, through
$R_H$, the whole construction is analytic in $H$. We speculate there is
a single \emph{Hurst-graded curvature operator} $\mathcal K(H,j)$,
analytic on $(H,j)\in(0,\infty)^2$, whose spectrum interpolates the
forward differences of $k\mapsto k^{2Hj/\cdots}$ and whose positive
cone is the common home of all three flows; the entire de
Bruijn--Newman / iterated-Brownian / random-matrix-edge complex is then
the bifurcation diagram of $\mathcal K$ as $H$ crosses the wall at $1$.
The Cauchy champion, the heat criticality, and the biharmonic collapse
are three points on one analytic family, and RH is the persistence of a
distinguished section in the positive cone. \emph{Peel-back note:} we do
not have $\mathcal K$ in general---but Remark~\ref{rem:chaosvar} now
gives the section $\mathcal K(\tfrac12,\,\cdot\,)$ in closed form,
spectrum $k^j-(k-1)^j$ for all real $j>0$, trace $N^j$, with its own
positivity wall at $j=1$. One axis of the dreamt operator is therefore
real, not dreamt; what remains dreamt is the $H$-direction and the
crossing of the wall.

%==============================================================
\section{What is proved, what is conjectured, what is dreamt}
\label{sec:ledger}
%==============================================================

An explicit ledger, so the speculation cannot be mistaken for the
content.

\begin{center}
\begin{tabular}{p{0.62\textwidth}l}
\hline
Statement & Status\\
\hline
Graded spectrum $k^j-(k-1)^j$ at $H=\tfrac12$, all $j$
(Thm.~\ref{thm:graded}) & \textbf{proved}\\
$\min$-eigenvalue $1$, $\tr=N^j$, $\det=\prod(k^j-(k-1)^j)$ &
\textbf{proved}\\
Wiener-chaos grading $R^{\Had j}=\tfrac1{j!}\mathrm{Cov}(:\!B^j\!:)$
\eqref{eq:chaos} & \textbf{proved (classical)}\\
Innovation factorization $R^{\Had j}=LD_jL^{\top}$, $j$-independent
eigenbasis, commuting tower (Lem.~\ref{lem:chol}, Cor.~\ref{cor:innov}) &
\textbf{proved}\\
Number operator $M^{(j)}=(\mathcal N+1)^j-\mathcal N^j$ generates the
tower (Prop.~\ref{prop:numberop}) & \textbf{proved}\\
Fifth fingerprint: all-ones factorization $\Leftrightarrow H=\tfrac12$
(Prop.~\ref{prop:fifth}) & \textbf{proved}\\
Sixth: tower is totally nonnegative, all minors via Cauchy--Binet,
Cauchy-only (Prop.~\ref{prop:TN}) & \textbf{proved}\\
Secular polynomial $p_j(z)=z^j-(z-1)^j$ has all zeros on $\Re=\tfrac12$;
critical line $\Rightarrow$ tower positivity (Prop.~\ref{prop:critline}) &
\textbf{proved}\\
Curvature spectral zeta $Z_2(s)=(1-2^{-s})\zeta(s)$; RH $=$ its critical
zeros on the line (Prop.~\ref{prop:czeta}) & \textbf{proved (identity);
RH open}\\
Real-$j$ tower, wall at chaos level $j=1$ (Rem.~\ref{rem:chaosvar}) &
\textbf{proved}\\
$\Psi_4$ rank $3$, signature $(1,2)$ (Prop.~\ref{prop:signature}) &
\textbf{proved}\\
All-orders optimality of $\gamma=1$ (Conj.~\ref{conj:phase}) &
numerically certified, $j\le4$, $N\le40$\\
$t^\ast(x)\asymp x^{3/2}$, i.e.\ $\Lambda(4)=\infty$
(Conj.~\ref{conj:lambda4}) & \textbf{refuted} ($\beta\approx0$, not
$x^{3/2}$; \S\ref{sec:gpu})\\
$\Lambda(4)<\infty\approx155=\tfrac{64}{27}\pi e^3$, saddle-detachment,
$x$-independent (Prop.~\ref{prop:lam4finite}) & \textbf{GPU + saddle},
to double precision ($x\le6000$)\\
$\Lambda(2m)<\infty$ all even orders, $t_*(m)$ closed form & saddle
detachment (\S\ref{sec:gpu})\\
$S1$: two Schur theorems, one phenomenon & speculation, plausible\\
$S2$: tower positivity $=$ real-rootedness & \textbf{false}; replaced by
the total-positivity dictionary (\S\ref{sec:dict})\\
Even-order coincidence Schur $=$ P\'olya--Schur (Thm.~\ref{thm:dict}) &
\textbf{proved}\\
$S3$: RH as critical tangency at $H=1$ \eqref{eq:S3} & speculation,
needs $S2$\\
$S4$: zero dynamics as Ricci flow / curvature healing & analogy\\
$S5$: one analytic operator $\mathcal K(H,j)$ & dream\\
\hline
\end{tabular}
\end{center}

The first block stands on its own and sharpens~\cite{FontesBE} from a
single eigenvalue identity to a graded law with an explicit chaos
meaning. The natural next theorem, and the one that would begin to
justify the ladder, is $S2$ for polynomials: to show that the Hankel /
Jensen determinants certifying $\LP$-membership of a real polynomial
are the leading minors of an associated Hadamard-power covariance, so
that real-rootedness becomes total positivity of a curvature tower. If
that holds, $S3$ acquires a precise form and the heat point becomes a
genuine criticality; if it fails, the structural picture of
\S\ref{sec:intro}--\ref{sec:ibm}---one Hurst line, a champion at
$\tfrac12$, a wall at $1$, a collapse at $2$---survives intact as the
part worth keeping.

%==============================================================
\section*{Acknowledgements}
%==============================================================
Claude was used as an interactive assistant for the computations,
drafting, and for generating and stress-testing the arguments, under
the author's direction. The author is responsible for all mathematical
content and any errors. The graded spectrum and the $(1,2)$ signature
were verified symbolically and in \texttt{numpy}/\texttt{scipy}.

%==============================================================
\begin{thebibliography}{99}

\bibitem{FontesBE} R.~Fontes, \emph{Bakry--\'Emery curvature of the
fractional Laplacian via fractional Brownian covariance}, 2026.

\bibitem{FontesIBM} R.~Fontes, \emph{Nonmonotone Laguerre--P\'olya
dynamics for the quartic de Bruijn flow} (iterated Brownian motion and
the Laguerre--P\'olya class), 2026.

\bibitem{FontesRMT} R.~Fontes, \emph{Quartic heat flow at the
random-matrix edge: speed, stability, and bulk lift-off}, 2026.

\bibitem{Schoenberg} I.~J.~Schoenberg, \emph{Metric spaces and positive
definite functions}, Trans.\ Amer.\ Math.\ Soc.\ \textbf{44} (1938),
522--536.

\bibitem{Mishura} Y.~Mishura, \emph{Stochastic Calculus for Fractional
Brownian Motion and Related Processes}, Lecture Notes in Math.\
\textbf{1929}, Springer, 2008.

\bibitem{RodgersTao} B.~Rodgers and T.~Tao, \emph{The de Bruijn--Newman
constant is non-negative}, Forum Math.\ Pi \textbf{8} (2020), e6.

\bibitem{Nualart} D.~Nualart, \emph{The Malliavin Calculus and Related
Topics}, 2nd ed., Springer, 2006.

\bibitem{ASW} M.~Aissen, I.~J.~Schoenberg, and A.~M.~Whitney, \emph{On the
generating functions of totally positive sequences I}, J.\ Analyse Math.\
\textbf{2} (1952), 93--103; A.~Edrei, \emph{On the generating functions of
totally positive sequences II}, J.\ Analyse Math.\ \textbf{2} (1952),
104--109.

\bibitem{PolyaSchur2} G.~P\'olya and J.~Schur, \emph{\"Uber zwei Arten von
Faktorenfolgen in der Theorie der algebraischen Gleichungen}, J.\ Reine
Angew.\ Math.\ \textbf{144} (1914), 89--113.

\bibitem{BorceaBranden2} J.~Borcea and P.~Br\"and\'en, \emph{P\'olya--Schur
master theorems for circular domains and their boundaries}, Ann.\ of
Math.\ (2) \textbf{170} (2009), 465--492.

\bibitem{CsordasVarga2} G.~Csordas, T.~S.~Norfolk, and R.~S.~Varga,
\emph{The Riemann hypothesis and the Tur\'an inequalities}, Trans.\ Amer.\
Math.\ Soc.\ \textbf{296} (1986), 521--541.

\bibitem{GORZ} M.~Griffin, K.~Ono, L.~Rolen, and D.~Zagier,
\emph{Jensen polynomials for the Riemann zeta function and other
sequences}, Proc.\ Natl.\ Acad.\ Sci.\ USA \textbf{116} (2019),
11103--11110.

\bibitem{BGKP} A.~Belton, D.~Guillot, A.~Khare, and M.~Putinar,
\emph{Matrix positivity preservers in fixed dimension I}, Adv.\ Math.\
\textbf{298} (2016), 325--368.

\end{thebibliography}

\end{document}
